\section{结论}

\subsection{研究结论}

本文围绕多机械臂搬运调度与模式选择问题展开研究。该问题不是机械臂运动学或轨迹控制问题，而是将机械臂视为可分配、可协同、可调度的作业资源，从运筹学角度研究任务分配、执行顺序、资源互斥、协同搬运和模式选择问题。

在主问题建模方面，本文将多类别物块搬运任务划分为单臂搬运任务和双臂协同搬运任务，并建立了包含机械臂资源互斥、中心安全区占用、任务顺序衔接和序贯目标规划的调度模型。模型首先最小化最大完工时间，其次在不破坏时间目标的基础上最小化系统总能耗，最后考虑机械臂负载均衡，从而体现了实际搬运系统中“先保证工期，再考虑能耗和均衡”的决策逻辑。

在主问题求解方面，本文采用 CP-SAT 求解器、融合剪枝的隐枚举法、启发式算法和热启动后的隐枚举法进行对比。结果表明，\textbf{CP-SAT 求解器、融合剪枝的隐枚举法和热启动后的隐枚举法在 16 种任务组合下均得到一致的最优结果}，说明自编精确搜索方法的求解结果是可靠的。启发式算法虽然计算速度最快，但在部分算例中不能完全达到最优解，因此更适合作为快速可行解生成方法或热启动初始解。

从计算效率看，融合剪枝的隐枚举法在任务数量较少时速度较快，但随着任务数量增加，搜索空间迅速扩大，计算时间明显上升。热启动后的隐枚举法利用启发式解提供较好的初始上界，在保持最优性的同时提高了剪枝效率，整体平均计算时间低于 CP-SAT 求解器和普通隐枚举法。因此，\textbf{热启动后的隐枚举法是主问题中综合表现较好的算法}。

在二臂与三臂模式选择方面，计算结果表明，\textbf{三臂系统平均缩短最大完工时间约 $31.24\%$，但平均能耗增加约 $4.46\%$}。这说明增加第三只机械臂的主要价值在于提高并行能力和缩短工期，而不是直接降低能耗。因此，\textbf{三臂系统并不是无条件优于二臂系统}，其推荐与否取决于决策者对时间收益和能耗代价的权衡。

在非线性延伸问题中，本文进一步引入速度倍率变量，将离散调度结果扩展到连续速度优化。问题一建立了速度--能耗非线性优化模型，并采用内点制约函数法结合 $0.618$ 法求解。KKT 验证结果表明，该问题在全部参数组下通过率达到 $99.85\%$，说明所得速度解整体满足一阶最优性条件，数值结果较为可靠。

问题二进一步考虑不同 seed 下任务空间分布变化，建立鲁棒速度模式选择模型，并采用逐次线性规划法求解。其 KKT 总通过率为 $64.65\%$，低于问题一。这说明鲁棒问题由于包含多个 seed 的时间约束和最坏能耗约束，活跃约束结构更加复杂，因此部分结果更适合解释为鲁棒可行解或数值近似解，而不能全部声称为严格最优解。但从工程意义上看，鲁棒模型能够反映不同任务位置扰动下二臂和三臂系统的稳健性差异，因此仍具有分析价值。

综上，本文的核心结论是：\textbf{多机械臂搬运系统的优化不能简单理解为“机械臂越多越好”或“速度越快越好”}。二臂系统在能耗控制方面具有优势，三臂系统在缩短工期和提高鲁棒性方面具有优势。\textbf{实际系统应根据任务结构、协同任务比例、截止时间要求、能耗代价和稳定性需求进行综合决策。}


\subsection{研究亮点}

本文的亮点并非简单增加求解方法，而是围绕同一决策链逐层推进：先完整描述多机械臂协同调度，再利用问题结构提高最优解搜索效率，最后把离散调度结果延伸到连续运行参数与鲁棒模式选择。

\begin{enumerate}
	\item \textbf{把任务分配、双臂组合和执行顺序放入同一个调度模型。}
普通任务分配只需要决定“每个任务由谁完成”，但本文还存在两类相互影响的决策：双臂任务需要选择哪两只机械臂协同，任务之间还需要确定先后顺序和开始时间。机械臂组合一旦改变，各臂的可用时间、后续空载移动和负载都会改变；任务顺序一旦改变，完工时间、转移能耗和中心区冲突也会随之改变，因此这些决策不能分开处理。

为此，本文把一种具体执行方案定义为一个“执行模式”：它明确某个任务由哪一只或哪两只机械臂完成，并给出相应的处理时间、能耗和中心区占用情况。模型为每个任务选择一个执行模式，同时安排所有已选模式的开始时间和先后关系，再通过机械臂互斥、双臂同步和中心区互斥约束排除冲突。这样，一个模型即可同时回答三个问题：任务由谁执行、双臂任务如何组合、所有任务按照什么顺序执行。它的好处是避免“先分配机械臂、再安排顺序”造成前一步看似合理、后一步却因等待、转移或中心区冲突而整体变差，使机械臂组合与任务顺序能够从全局目标出发共同优化。最后再按最大完工时间、总能耗和负载均衡的优先级依次优化，使整数规划和目标规划直接服务于实际调度决策。

	\item \textbf{利用前瞻预测和启发式上界，更快获得并证明最优解。}
如果直接完整枚举，需要同时遍历任务顺序、机械臂分配方式和双臂协同组合，搜索空间会随任务数量迅速膨胀。本文的关键改进不是等一个分支已经展开很深后再判断，而是在选择“下一个任务及其执行模式”时就提前预测该选择的后续潜力。

这种预测分为两层。第一层计算候选任务真实插入当前调度后的最早开始时间、完成时间、空载转移能耗和负载变化，判断这一步本身是否合理；第二层暂不构造具体后续顺序，而是把所有剩余任务放松，估计该候选继续完成全部任务时不可能突破的最好结果。预测同时考虑单任务最早完成、机械臂容量、必须使用某只机械臂的任务、中心区串行占用和剩余最低能耗。候选进入递归搜索之前，程序先用这一前瞻下界排序并筛除明显没有希望的选择；进入子状态后，再用完整下界和状态支配规则复检。与只在搜索节点上被动剪枝相比，这相当于在“是否值得走入该分支”之前增加了一次安全的预判。

这里的“预测”不是猜测未来任务的具体排法，而是计算一个偏乐观的理论下界。真实后续调度只可能等于或差于该下界，因此只有当这个最好可能结果都无法优于当前完整解时，分支才会被剪去，不会因此误删潜在最优解。这种前瞻评价使算法既能提前减少候选数量，又保留精确搜索的最优性保证，是融合剪枝隐枚举法区别于直接枚举的重要设计。

热启动进一步放大了前瞻预测的作用。启发式算法先快速给出较好的完整可行解，其字典序目标值成为初始上界，于是前瞻下界从搜索开始就有明确的比较对象。启发式负责“尽早给出好上界”，隐枚举负责“利用下界排除不可能更优的分支并完成最优性证明”，形成“先快速找到较好解，再围绕较好解进行精确搜索”的过程。隐枚举仍从空调度开始，不改变可行域、安全下界或状态支配规则，所以启发式解不再只是近似答案，而被转化为精确搜索的剪枝依据；搜索完成后，最优性仍由完整分支定界保证。由此，热启动隐枚举不仅用于改进启发式解，也能够反过来判定启发式是否已经达到最优目标值。

在本文算例中，热启动后的隐枚举方法平均快于 CP-SAT，正是因为前瞻预测、问题专用下界和紧初始上界共同减少了实际展开的分支。CP-SAT 则需要将同一结构转换为通用可选区间、回路和逻辑约束，并分阶段处理三级目标。实验中热启动隐枚举平均用时为 $6.0634$ s，低于 CP-SAT 的 $7.2631$ s，且 16 种任务组合的最优目标值完全一致。这里的优势来自本文问题结构与算例规模，不能外推为自编算法在任意规模下都快于通用求解器。

\textbf{启发式最优性}采用内部证明与外部对照两种方式判断。内部证明来自热启动隐枚举：程序分别保留启发式初始目标值和隐枚举完成后的最优目标值；若两者一致，说明完整搜索没有发现任何更优分支，启发式已经达到该算例的字典序最优目标；若最终目标值更优，则说明启发式只是近似解，而隐枚举完成了改进。外部对照则使用 CP-SAT 已证明的最优结果，按照最大完工时间、总能耗和负载差的优先级依次比较。两种验证相互补充：隐枚举给出自编算法内部的最优性证明，CP-SAT 提供独立求解器的交叉检查。实验据此统计最优命中率、目标差距和计算时间，二臂场景有 $44/48$ 组启发式解达到最优，三臂场景有 $23/48$ 组达到最优。

	\item \textbf{把离散调度结果继续推进为调速与鲁棒选型决策。}
主问题能够确定任务由哪只机械臂执行、双臂任务采用哪种组合以及任务按什么顺序执行，但这一阶段默认运行速度给定，只回答了“任务如何安排”，还不能回答“既定安排下以多大速度运行更合适”。如果速度过低，系统可能无法满足截止时间；如果一味提高速度，又会增加驱动能耗。因此，本文没有停留在离散调度结果的比较，而是把最优调度得到的时间和能耗结构继续作为连续优化的输入。

在第一层延伸中，模型引入单臂任务与双臂任务的速度倍率，在满足截止时间的条件下优化运行能耗，实现从任务层决策到运行参数决策的衔接。主问题决定“谁来做、先做什么”，非线性模型进一步决定“按什么速度做”。这样既保留了已优化的调度结构，又能单独识别速度变化带来的时间--能耗影响。在第二层延伸中，考虑同一任务规模在不同 seed 下具有不同空间分布，不再针对每个 seed 分别选择一组有利速度，而是寻找能够同时适应多组位置扰动的统一速度设置，并控制最坏情形能耗。统一速度方案的好处是避免对某一个任务位置样本过度优化，使所得参数在任务位置变化后仍可直接使用。这样，研究从单一算例下的最优调速继续推进到不确定任务位置下的鲁棒模式选择。

这三层问题使用不同工具，但输入输出关系是连续的：整数规划和隐枚举先得到离散调度方案，内点制约函数法结合 $0.618$ 法求解确定性速度--能耗问题，逐次线性规划再处理多 seed 的鲁棒约束。采用分层求解而不是把任务组合、执行顺序和连续速度一次性放入庞大的混合整数非线性模型，能够降低单次求解复杂度，并使每一阶段的目标、约束和结果来源都更容易解释。验证方式也与问题层次对应：主问题用 CP-SAT 交叉验证最优目标，确定性非线性问题用 KKT 条件检查一阶最优性，鲁棒问题则根据 KKT 检查区分严格一阶最优解与数值可行解。因此，这一部分的亮点不是简单增加两个延伸问题，而是在兼顾可求解性和可解释性的前提下，形成“先安排任务、再优化运行参数、最后检验扰动下稳定性”的完整决策链。

该决策链最终服务于二臂/三臂系统选型。实验表明，三臂系统平均缩短最大完工时间约 $31.24\%$，同时增加能耗约 $4.46\%$；速度优化和鲁棒分析又进一步说明，这一取舍会受到截止时间与任务位置扰动的影响。因此，项目没有停留在“机械臂越多越好”或“速度越快越好”的笼统判断，而是把系统选择转化为可计算、可验证的时间收益--能耗代价权衡。
 
\end{enumerate}

\subsection{心得体会}

通过本次大作业，我首先从代码实现层面加深了对课内知识的理解。课堂上学习整数规划、目标规划、隐枚举法、启发式算法、非线性规划和 KKT 条件时，更多是从定义、公式和例题角度理解方法；而在本项目中，这些内容都需要落实到具体变量、约束、目标函数和程序逻辑中。例如，目标规划不只是写出多个目标，而是要在程序中体现目标的优先级；隐枚举法也不是简单地枚举所有方案，而是要设计合理的下界和剪枝规则；KKT 条件也不只是理论公式，而是可以用于检查非线性优化结果是否满足一阶最优性。通过把课内方法转化为可运行的程序，我对这些方法的理解更加具体。

其次，本次作业让我认识到，运筹学建模必须从实际问题结构出发。最初看多机械臂搬运问题时，很容易把它理解为普通任务分配问题，认为只要决定每个任务由哪只机械臂执行即可。但在实际建模过程中可以发现，任务顺序、机械臂当前位置、双臂协同、资源互斥、中心安全区占用和目标优先级都会影响最终结果。因此，一个实际问题往往不能直接套用单一标准模型，而是需要先分析问题中的关键矛盾，再选择合适的数学工具进行抽象。这个过程使我更加理解了“建模”本身的重要性。

最后，本次作业也让我体会到，求解结果必须经过对比和验证。一个算法能够输出结果，并不代表结果一定可靠。对于主问题，需要将自编算法与 CP-SAT 求解器进行对照，判断其是否达到标准求解器的结果；对于非线性延伸问题，需要用 KKT 条件检查可行性、互补松弛和平稳性，判断数值解是否具有一阶最优性。通过这种验证过程，才能说明模型和算法不是只停留在形式上，而是具有一定可信度。

从后续扩展角度看，本文仍有进一步改进空间。当前模型主要从任务层调度角度进行抽象，还没有深入考虑真实机械臂的避碰轨迹、加减速过程和动力学约束；任务集合也是静态给定的，没有考虑动态任务到达或临时任务插入；鲁棒问题目前采用逐次线性规划得到数值近似解，后续还可以引入更稳定的非线性求解器或更严格的全局优化方法进行验证。如果进一步扩展，可以使模型更加接近真实工业搬运场景。

\section{思考}

\subsection{主问题是否可以理解为一部分动态规划？}

本文主问题可以理解为具有一定动态规划思想，但不能严格称为标准动态规划问题。

从相似性来看，主问题确实存在“分阶段决策”的特征。调度过程可以看作逐步选择任务的过程：在每一步中，需要根据当前系统状态选择一个尚未完成的任务，并决定由哪只机械臂或哪两只机械臂执行。任务执行完成后，系统状态发生变化，例如已完成任务集合增加，机械臂当前位置改变，机械臂可用时间更新，中心安全区释放时间也可能发生变化。后续任务的安排又会受到这些状态变化的影响。因此，从“状态—决策—状态转移”的角度看，本文主问题与动态规划具有一定共同点。

如果将其写成动态规划思想下的状态，状态至少需要包含
\[
\Omega=
(\text{已完成任务集合},
\text{各机械臂当前位置},
\text{各机械臂可用时间},
\text{中心安全区状态},
\text{累计能耗},
\text{负载情况})。
\]
每一步决策则可以表示为：从未完成任务中选择一个任务，并选择对应的执行模式。对于单臂任务，需要选择一只机械臂；对于双臂协同任务，需要选择两只机械臂。决策完成后，系统进入新的状态。从这个意义上说，本文主问题确实包含动态规划的阶段决策思想。

但是，本文主问题并不适合直接采用标准动态规划求解。原因在于状态维度过高，状态数量会随着任务数量和机械臂数量迅速增加。仅“已完成任务集合”这一项就会带来指数级状态数量；再加上机械臂位置、可用时间、中心安全区占用、能耗和负载等连续或半连续信息，状态空间会变得非常庞大。如果按照标准动态规划建立状态表并逐层递推，计算量会难以接受。

此外，标准动态规划通常要求状态转移关系相对明确，并且每个状态的价值函数能够被有效存储和复用。而本文主问题中，任务执行时间和能耗不仅取决于当前任务，还取决于前序任务位置、机械臂组合、双臂同步等待和中心安全区冲突。因此，不同路径到达的状态即使已完成任务集合相同，也可能由于机械臂位置和时间不同而具有完全不同的后续代价。这使得状态合并和价值函数递推都比较困难。

因此，本文主问题更准确的理解是：它具有动态规划的阶段决策思想，但实际求解上更接近组合调度中的分支搜索和隐枚举方法。本文采用的隐枚举法也是逐步构造调度方案，但它不是建立完整动态规划状态表，而是在搜索树中结合下界和剪枝规则排除不可能成为最优解的分支。因此，主问题不能严格称为动态规划问题，但可以说它吸收了动态规划中“逐阶段决策、状态随决策变化”的思想。

\subsection{匈牙利算法为何不能使用？}

匈牙利算法是求解标准指派问题的经典方法，尤其适用于一类特殊的 0--1 整数规划问题。在标准指派问题中，通常有若干个执行者和若干个任务，每个执行者分配一个任务，每个任务也只由一个执行者完成，并且执行者完成任务的成本在求解之前已经给定。其典型模型为
\[
\min \sum_i\sum_j c_{ij}x_{ij},
\]
其中 $c_{ij}$ 表示执行者 $i$ 完成任务 $j$ 的固定成本，$x_{ij}$ 表示是否将任务 $j$ 分配给执行者 $i$。在这种情况下，问题的核心是静态的一对一匹配，匈牙利算法能够高效求解。

本文主问题虽然也包含 0--1 决策变量，但它并不是标准指派问题，因此不能直接使用匈牙利算法。首先，本文中的任务并不都是一对一执行。Type1 和 Type2 可以由单只机械臂完成，但 Type3 和 Type4 需要两只机械臂协同完成。也就是说，一个任务可能需要两个执行者同时参与，这已经突破了标准指派问题中“一个任务对应一个执行者”的基本假设。

其次，本文中的成本不是固定成本矩阵。某只机械臂执行某个任务所需的时间和能耗，不仅取决于任务本身，还取决于该机械臂上一项任务的位置、当前可用时间、是否需要等待另一只机械臂协同，以及中心安全区是否被占用。因此，成本会随着任务顺序和系统状态变化，而不是求解前就固定好的 $c_{ij}$。没有固定成本矩阵，匈牙利算法就失去了直接适用的前提。

再次，匈牙利算法只能解决“谁分配给谁”的问题，不能解决“任务按什么顺序执行”的问题。而本文主问题中，任务顺序是影响最大完工时间和总能耗的关键因素。同样的任务分配，如果执行顺序不同，机械臂的空载移动距离、等待时间、中心区冲突和最终完工时间都会不同。因此，只做一次静态分配无法描述完整调度过程。

最后，本文还包含机械臂资源互斥、中心安全区互斥、双臂协同同步和序贯目标规划等约束。机械臂不能同时执行多个任务，中心安全区不能被多个任务同时占用，双臂任务需要两只机械臂协调执行，目标函数还要按照最大完工时间、总能耗和负载均衡的优先级依次优化。这些内容都超出了匈牙利算法的适用范围。

因此，匈牙利算法虽然是求解标准 0--1 指派问题的有效方法，但它适用于固定成本、静态、一对一匹配问题。本文主问题同时包含任务分配、任务排序、双臂协同、资源互斥、中心安全区冲突和多目标优先级，因此不能直接使用匈牙利算法。最多只能在极度简化的局部场景中，例如某一时刻将若干单臂任务临时分配给空闲机械臂时，将其作为局部分配工具参考，而不能作为完整主问题的求解方法。
```
